\begin{definition} Класс \np состоит в точности из тех языков $A$, для которых существует детерминированная машина Тьюринга двух аргументов $V(x, s)$, работающая по времени за $poly(|x|)$ такая, что
$$
\forall x \in \Sigma^* \ x \in A \Longleftrightarrow \exists s: V(x, s) = 1.
$$
Тогда $V$~--- верификатор, а $s$~--- сертификат для данного $x$.
\end{definition}

\begin{definition} Классом $\ntime{T(n)}$ называется класс языков, которые распознаются за время $O(T(n))$ недетерминированной машиной Тьюринга. Иными словами, время работы машины на любом слове длины $n$ не превосходит некоторой константы, умноженной на $T(n)$.
\end{definition}

\begin{definition} $\np = \bigcup\limits_{c=1}^{\infty}\ntime{n^c}$
\end{definition}

\begin{theorem}\label{th:eqvivNP}
Эти определения эквивалентны.
\end{theorem}
\begin{proof}

Пусть $A$ распозанется недетерминированной машиной Тьюринга, которая работает $O(n^c)$ шагов. Тогда в качестве сертификата можно взять последовательность значений функции перехода, а верификатор будет моделировать работу машины $M$, используя данные сертификата для выбора одной из ветвей алгоритма. Можно действовать и по-другому: в качестве сертификата взять всё
вычисление, а верификатором проверять его соответствие программе.

Пусть для $A$ верно сертификатное определение. Тогда сертификат не может быть длиннее, чем время работы $V$, т.е. чем некоторый полином $p(|x|)$. Недетерминированная машина будет работать следующим образом: сначала она недетерминированно напишет сертификат $s$ длины не больше $p(|x|)$, а затем запустит $V(x,s)$. 
\end{proof}